\section{차별성과 결론}

\subsection{창의성 및 차별성}

C-BAT의 차별점은 새 알고리즘이 아니라, 시간별 생산계획을 입력으로 받는 해석 가능한 구조에서 하루 경로의 분포를 만들고 그 분포를 요금이 정해지는 방식에 맞춰 쓴 조합에 있다. 표~\ref{tab:ch6-compare}는 이 과제에서 쓸 수 있는 기존 접근과 C-BAT를 다섯 가지 기능으로 비교한다. 표시는 각 접근의 일반적인 형태를 기준으로 하며 기존 BAT의 월 최대 경보는 경로를 만들 수 있지만 이 보고서에서 적용하지 않았다는 뜻이다.

\begin{table}[h]
\centering
\caption{기존 접근과 C-BAT의 기능 비교(○ 제공, △ 부분 제공 또는 미적용, × 없음).}
\label{tab:ch6-compare}
\setlength{\tabcolsep}{4pt}
\small
\begin{tabular}{lccccc}
\toprule
접근 & 시간별 계획 & 예측분포·경로 & 피크 위험확률 & 월 최대 경보 & 해석 가능성 \\
\midrule
유사일·B0 기준일 & × & × & × & × & ○ \\
GAM·평활 가법모형 \citep{hyndman2010peak} & △ & △ & △ & × & ○ \\
LightGBM \citep{ke2017lightgbm} & ○ & △ & △ & × & △ \\
기존 BAT & ○ & ○ & ○ & △ & ○ \\
C-BAT & ○ & ○ & ○ & ○ & ○ \\
\bottomrule
\end{tabular}
\end{table}

전체 방법의 효과는 합산 평가로 확인했다. 개발 구간과 9월을 합친 피크 평가 65일의 MAE는 C-BAT 8.93 kW, 기존 BAT 13.09 kW, M2 14.54 kW였다. 기존 BAT 대비 31.8\%, M2 대비 38.6\% 개선이며, 날짜 단위 짝지은 bootstrap 95\% 구간에서도 두 차이가 0을 포함하지 않는다. 개발 구간 고피크일 13일에서는 C-BAT 4.37 kW, 기존 BAT 12.47 kW, 튜닝 LightGBM 18.27 kW였다.

설계 요소별 절제는 선택 근거로 해석한다. 가동 상태별 잡음은 비가동일과 가동일에 서로 다른 변동을 주고, 고정 커널은 학습해야 할 형상 모수 9개를 줄인다. 학습형 커널에서 공통 잡음으로 바꾸면 저장된 고피크 MAE가 4.38 kW에서 7.80 kW로 커졌지만, 두 비교 모형은 split-$\hat R$ 기준을 통과하지 못했다. 따라서 이 차이를 전체 개선 8.10 kW 가운데 잡음 설계가 기여한 3.42 kW라고 분해하지 않는다. 하루 경로는 요금시간대 최댓값의 분포와 초과 비용을 같은 표본에서 계산하게 한다. 이 경보는 7--8월의 유일한 월 최대 갱신일인 7월 19일을 잡았고 B0\_kind는 이날을 놓쳤다. 사건 한 건의 포착은 운용 사례이며 절감 효과의 검증은 아니다.

평가 절차도 차별점이다. 복사일 122일을 찾아 누수 규모를 정량화한 뒤 시간순 검증을 설계했고 \citep{kapoor2023leakage}, 기준일 규칙·잡음 변형·경보 규칙은 결과를 보기 전에 채택 기준을 기록한 뒤 평가했다. 모형의 반사실 일정 이동 절감액은 관측 유사일 검증에서 지지되지 않아 주장하지 않았다. 에너지 인지 스케줄링 연구가 설비 전력을 알려진 값으로 두는 것 \citep{fang2011peakpower,shrouf2014tou}과 달리, 이 보고서는 모형 출력을 관측으로 확인한 범위에서만 현장 제안에 썼다.

\subsection{한계와 결론}

평가 범위는 한 공장의 2021년 자료다. 주 성능은 개발+9월의 67일·6,358슬롯, 피크 65일에서 비교했다. 모델 선택과 절제는 개발 구간의 달력상 56일 가운데 유효 53일, 피크 51일과 고피크일 13일에서 수행했다. 개발 구간은 여러 모델 비교에 반복해서 쓰였으므로 선택과 조건별 비교는 탐색적이다. 합산 비교에는 B0\_kind와 튜닝 LightGBM의 9월 결과가 포함되지 않아 모든 모델에 대한 최우수 성능으로 표현하지 않는다. 다른 계절과 다른 공장으로 일반화되는지도 확인하지 않았다.

방법에도 한계가 남아 있다. C-BAT는 고피크가 아닌 평범한 가동일의 피크를 높게 예측하는 경향이 있고 이것이 7월 초 경보의 오경보 9건으로 이어졌다. 개발 구간 전체 일 피크 MAE는 단순 기준선 B0\_kind가 6.67 kW로 더 작았다. 슬롯 90\% 구간 적중률은 개발 구간 0.990, 9월 0.984로 명목값보다 높아 구간이 넓었다. 합산 슬롯 MAE는 기존 BAT보다 0.61 kW 컸다.

데이터 제약도 있다. 사전 생산계획이 없어 생산실적을 계획의 대용으로 썼다. 설비별 전력이나 설비 동시 가동 정보가 없어 피크를 만드는 설비를 특정할 수 없다. 완전 복사일 115일과 편집 복사일 7일, 총 122일을 제외해 학습 자료가 줄었다.

기간별 성능도 함께 보고했다. 9월 단독에서는 M2가 슬롯 MAE 6.60 kW, 일 피크 MAE 5.53 kW로 C-BAT(7.18 kW, 8.93 kW)보다 좋았다. 이 결과는 합산 평가에서 확인한 피크 개선과 함께 읽어야 한다. 9월에는 고피크일이 하루도 없어 해당 조건의 성능을 새로 시험할 수 없었고 이전 개발 때 열린 적이 있어 완전히 독립된 홀드아웃도 아니다. 월 최대 갱신 사례는 1건뿐이며, 주 경보 규칙의 가정대로 점검 10회에 회당 3만 원을 지출하면 30만 원이다. 그 달 완벽 조치 금액 상한 133,120원보다 커서 당월 순편익이 양수라는 근거는 없다. 이후 래칫 영향과 실제 조치 효과는 별도 검증이 필요하다.

향후 과제는 세 가지다. 설비별 전력과 가동 기록을 받아 기동 부하와 동시 가동을 모형에 넣는다. 경보일을 조치일과 비조치일로 나누는 현장 A/B 운영으로 조치의 실제 효과를 측정한다. 여러 계절과 여러 번의 월 최대 갱신을 포함하는 긴 기간으로 경보 규칙을 다시 평가한다.

이 보고서는 시간별 생산계획으로 15분 전력과 일 피크의 확률분포를 예측하는 C-BAT를 제시했다. 개발+9월 합산 평가에서 피크 MAE는 8.93 kW로 기존 BAT 13.09 kW와 M2 14.54 kW보다 작았고, 개발 구간 고피크일 MAE는 4.37 kW였다. 같은 경로로 만든 월 최대 갱신 경보는 유일한 갱신일을 포착했다. 보고할 성과는 이 예측 개선과 경보 사례이며, 133,120원은 완벽 조치를 가정한 금액 상한이다. 생산실적을 사전 계획으로 간주한 입력 가정, 평범한 가동일 과대예측, 점검 비용과 조치 효과를 현장 검증에서 확인해야 한다.
